\section{Data}\label{sec:data}
\paragraph{Source.}
Measurements were retrieved from the ChEMBL REST API on 10 October 2026; the API reported release ChEMBL~37 (released 1 May 2026)~\cite{zdrazil2024chembl}. We queried 59 single-protein targets (Appendix~\ref{app:repro}; identifiers checked against ChEMBL preferred names). Twenty-nine were chosen for the first version of this work; the further thirty were added afterwards by taking the human single-protein targets with the most binding $K_i$/$\mathrm{IC}_{50}$ activities carrying a pChEMBL value that were not already in the set (\texttt{scripts/}\allowbreak\texttt{13\_discover\_targets.py}; 242 of 5{,}869 targets have at least 1{,}500 such activities). We kept activities with standard type $K_i$ or $\mathrm{IC}_{50}$, relation `$=$', a pChEMBL value, binding assay type (B), no data-validity comment, no potential-duplicate flag, a document year, and an assay confidence score $\ge 8$. A score of 8 admits assays assigned to a homologous single protein, so assays on orthologues from other species may be included; we did not filter by organism or by mutant variant. Censored values (relations `$<$', `$>$') are dropped, which truncates the potency range at both ends.

\paragraph{Standardisation and aggregation.}
Structures were standardised with RDKit (largest fragment, neutralisation, canonical SMILES). Molecules are identified by the connectivity block of the InChIKey, so stereoisomers are merged (a stereo-aware sensitivity analysis is given in Section~\ref{sec:ceiling}). Within a document, values for a molecule are averaged, so that every remaining replicate is \emph{between} documents; note that one document can contain several assays (e.g.\ different enzyme variants or assay conditions), which this averaging mixes. A molecule's label is the mean over its documents; we also record the number of documents $n_i$ and the year of first report. A \emph{task} is a (target, measurement type) pair; $K_i$ and $\mathrm{IC}_{50}$ are never pooled. This gives 118 tasks.

\paragraph{Task selection.}
A task is kept if at least 100 molecules have values from $\ge2$ documents. Forty-one tasks pass after decontamination, 17 from the original 29 targets and 24 from the 30 added later (Table~\ref{tab:ceil}). They contain 1{,}000--9{,}567 molecules (median 4{,}778; 210{,}519 in total). Only 1.8--18.4\% of molecules per task (median 3.6\%) have replicates; after trimming (Section~\ref{sec:ceiling}) the floors are estimated from 8{,}916 replicated molecules. The 24 added tasks were analysed with the same pipeline and thresholds and serve as a replication set for the model comparison: they cover target classes the original set did not, among them a bromodomain, a deacetylase, two ion channels, a nuclear receptor and an isocitrate dehydrogenase.

\paragraph{Decontamination.}
We built a registry of molecules from 11 MoleculeNet and 40 TDC~\cite{huang2021tdc} datasets (202{,}019 unique structures) and removed, from the 41 tasks, molecules matching the registry by full InChIKey (5{,}393 molecules), additionally by InChIKey connectivity block (1{,}175), and near analogues with ECFP4 Tanimoto similarity $>0.8$ to any registry molecule (3{,}121): 9{,}689 of 220{,}208 molecules (4.4\%) in total. Replicated molecules were removed more often: 1{,}541 of 10{,}561 (15\%), since well-studied compounds are also the ones that appear in public benchmarks. Decontamination therefore changes the population from which the floor is estimated, and on the first 29 targets it reduced the number of tasks passing the threshold from 21 to 17. Polaris datasets could not be downloaded without authentication and are not included. Pretraining corpora of the foundation models (PubChem and ZINC subsets) were not screened; a filter removing documents published before the models' release left no task above the threshold and was therefore not applied. Unless stated otherwise, all results use the decontaminated data; the same models run on the uncleaned data are compared in Section~\ref{sec:results}.

\paragraph{Replicate statistics.}
An audit of the replicated pairs on the original 17 tasks before decontamination (\texttt{results/audit/}), in which pairs rather than molecules are trimmed, shows the following. The pooled between-document standard deviation depends strongly on how large disagreements are treated: its median over tasks is 0.65 log units without trimming, 0.57 when pairs differing by more than 3 log units are excluded and 0.45 when those differing by more than 2 are excluded. The disagreement is only weakly related to the time between the two publications (median Spearman $\rho=0.08$). Replicated molecules are not a random sample: on the decontaminated data they are on average 0.4--1.0 log units more potent than molecules reported only once in 7 of the 17 tasks.

\paragraph{Provenance of the replicates.} Whether two documents are genuinely independent, whether one re-publishes the other, and what assay conditions each used all bear on the floor, and none of them follows from the activity records alone. Section~\ref{sec:provenance} takes them up with document metadata: the first two can be measured, the third largely cannot. Species and mutant constructs we still do not check.
